\subsection{Zeit-Anzeige-Logik der WordClock}
Die zentrale Funktion der Software besteht in der Umwandlung der numerischen Zeit in eine sprachbasierte Darstellung. 
Dabei werden Minutenintervalle interpretiert und die entsprechende Stunde abhängig vom Kontext angepasst. 
Die Ausgabe erfolgt durch gezielte Aktivierung von \gls{led}-Gruppen, welche Wörter darstellen. 
Der Ablauf der Anzeigeerzeugung ist in Abbildung~\ref{fig:WordClock11} dargestellt.

\begin{figure}[H]
\centering
\begin{tikzpicture}[
    node distance=1.0cm,
    startstop/.style={rectangle, rounded corners, minimum width=4cm, draw=black},
    process/.style={rectangle, minimum width=4.8cm, draw=black},
    decision/.style={diamond, aspect=2.3, draw=black},
    io/.style={trapezium, trapezium left angle=70, trapezium right angle=110, minimum width=4.6cm, minimum height=0.8cm, text centered, draw=black},
    arrow/.style={->, thick}
]

\node (start) [startstop] {Anzeige starten};
\node (clear) [process, below=of start] {LEDs zurücksetzen};
\node (power) [decision, below=of clear] {LEDs aktiv?};
\node (always) [io, below=of power] {Feste Wörter anzeigen};
\node (calc) [process, below=of always] {Anzuzeigende Zeit ermitteln};
\node (minute) [io, below=of calc] {Minutenwort anzeigen};
\node (hour) [io, below=of minute] {Stundenwort anzeigen};
\node (weather) [io, below=of hour] {Sekunden anzeigen};
\node (end) [startstop, below=of weather] {Anzeige fertig};

\draw [arrow] (start) -- (clear);
\draw [arrow] (clear) -- (power);
\draw [arrow] (power) -- node[right]{Ja} (always);
\draw [arrow] (always) -- (calc);
\draw [arrow] (calc) -- (minute);
\draw [arrow] (minute) -- (hour);
\draw [arrow] (hour) -- (weather);
\draw [arrow] (weather) -- (end);

\draw [arrow] (power.west) -- ++(-2,0) node[above]{Nein} |- (end.west);

\end{tikzpicture}
\caption{Schematischer Programmablaufplan zur Erzeugung der Zeit-Anzeige-Logik der WordClock.}
\label{fig:WordClock11}
\end{figure}